\section{Notation and reading index}
\label{app:notation-index}

The original DRR field and graphic are \(F_*,\Gamma_{I4}\);
the canonical endpoint is \(H_0,\Gamma\); the
neighboring organizing family is \(H_c,\Gamma_c\). The full quadratic germ is
\(\mathcal Q_2\), while \(\theta=(A,C,\mu_1,\mu_2,\mu_3)\) denotes the
five-parameter transverse slice. In the positive-infinity chart,
\(v=X/Y\) and \(z=1/Y\); \(h\) denotes a central height only where that
local section is fixed. The exact return is written
\(\mathcal P_\theta(x)=m_\theta\Pi_\theta(x)\), where \(m_\theta\) is the
finite through multiplier and \(\Pi_\theta\) is the completed divided
transition; \(F_\theta=x-\mathcal P_\theta\) is the displacement.

The local finite saddle--node through multiplier is
\(m_{\rm fin}=e^{-T_f}\) in \eqref{eq:finite-multiplier}, written \(m\)
in the completed return \eqref{eq:return}. In the fixed-\(c\) lower construction,
\(\Lambda_\rho\) is the complete external gain, including the two tail
gains and the finite connection, as defined in \eqref{p4:outer-data}.
The action-model target is \(m_\eps=2\eps^{5/2}\); the two-coefficient
calibration \eqref{p4:calibration} imposes \(\Lambda_\rho=m_\eps\).

In the positive-infinity chart, the polynomials are
\[
 P(v)=\mu _1+\mu _2v-v^3,\qquad
 R(v)=\mu _3+v+v^2
\]
as in \eqref{eq:polynomials}.  The symbol \(\gamma\) denotes the small
root of the transverse-rate equation \(R(\gamma)=0\), equivalently the
root selected in \eqref{eq:center-parameters}; it is not, in general, a
root of the equatorial drift \(P\).  The two symbols
\(\alpha _0,\alpha _1\) are the value and first derivative of that drift,
\[
 \alpha _0=P(\gamma),\qquad \alpha _1=P'(\gamma),
\]
not transverse coefficients.  The ratios \(p,\pi,\mathfrak q,U,G\) are
local coefficient dictionaries introduced at their first use in the
coefficient-regime and degree-loss sections.

The scale \(\tauc=t^2\) is reserved for the cubic coefficient chart in
\eqref{deg:cubic-chart}.  An unadorned \(\tau\) is local and must be read
from its defining display: in the coefficient charts it denotes the stated
small coefficient scale (for example \(p^2\) or \(gt_c/a\)), whereas in a
variation integral or a reparametrized orbit it can be a dummy or clock
variable.  No global identification of these uses is intended.  The source
variable and the physical section coordinate are not identified unless the
displayed affine adapter says so.

Finally, \(\Pi_\theta\) in \eqref{eq:return} is the completed, divided
transition: the infinity map composed with the regular section adapters and
the finite-side closing factor after the factor \(m\) has been extracted.
It is the positive transition appearing in the exact displacement
\(F=x-m\Pi_\theta(x)\); it is not an unrestricted Poincare map outside the
chosen physical sections.

The coefficient dictionaries define their ratios before use. The common
raw receiver and capture retain the original physical coordinate throughout;
positive affine changes transfer the derivative signs used for counting.
